A focused related-work chapter for the specific contributions of this
study: glitch-free dual-camera switching on a shared MIPI-CSI lane,
per-iteration measurement of switch cost on MCU-class hardware, and the
platform-level design choices that made those measurements possible.
Complements the broader survey in \url{relatedwork.md}, which covers
nano-UAV autonomy, YOLO-family detectors, and MicroPython-in-robotics
context; this chapter is narrower and deeper on the switching problem.

\begin{center}\rule{0.5\linewidth}{0.5pt}\end{center}

\section{1. Hardware multiplexing of multiple cameras onto one CSI
receiver}\label{hardware-multiplexing-of-multiple-cameras-onto-one-csi-receiver}

\subsection{1.1 Analogue-switch MUX (our
topology)}\label{analogue-switch-mux-our-topology}

Multiplexing two camera modules onto a single MIPI-CSI2 lane via an
analogue switch is the default approach on small-form-factor evaluation
boards where the SoC has only one CSI receiver and board area is scarce.
The TS5MP645 family from Texas Instruments is a representative device;
in the NXP i.MX ecosystem, the same pattern appears on the i.MX8M Mini
where the OV5640 community thread (1064767) documents two-camera
operation through a single CSI controller with GPIO-selected mux. The
analogue-switch approach is attractive because the MUX is cheap,
passive, and transparent to the D-PHY receiver once the line settles,
but it carries two known limitations:

\begin{itemize}
\tightlist
\item
  \textbf{No virtual-channel separation.} Both sensors' MIPI short
  packets (FS, LS) arrive on the same virtual channel (VC=0) from the
  receiver's perspective; there is no hardware disambiguation between
  them. This forces the platform to ensure only one sensor is driving
  the line at a time, which is exactly what the GPIO MUX guarantees.
\item
  \textbf{Switch timing is software's responsibility.} The receiver has
  no notion of ``sensor change''; if the MUX flips mid-line, the DMA
  buffer being written at that moment silently contains pixels from both
  sensors. Our measurement of this failure mode is in
  \url{cam_switch.md} §``Visual evidence --- seam before the fix'' and
  is the ground truth that motivates flip-on-EOF.
\end{itemize}

The TI E2E thread on the TS5MP645 MIPI-CSI2 MUX (874026) identifies
D-PHY re-training after each switch as an optional but expensive
mitigation; we instead keep the MUX transition inside the MIPI blanking
window so no re-training is needed.

\subsection{1.2 Virtual-channel-separated dual-stream
receivers}\label{virtual-channel-separated-dual-stream-receivers}

Hardware that supports MIPI-CSI2 virtual channels natively --- e.g.
systems using MIPI-bridge chips such as the Lontium LT9611 or Toshiba
TC358748 --- can run two sensors concurrently on separate VCs and
demultiplex at the receiver side. On such a topology the ``switch cost''
problem disappears entirely because both streams are live. NXP's
MIPI-CSI2 peripheral on the RT1170/1176 supports four virtual channels
on the receive side, but the OV5640 cannot natively remap its output VC,
and we do not have a bridge chip on the SentAI board. This positions our
measurements as a lower-bound study: any platform that upgrades to a
VC-split topology should expect \emph{better} dual- camera throughput
than we report, not worse.

\subsection{1.3 Dual CSI controllers}\label{dual-csi-controllers}

SoCs with multiple independent CSI receivers (e.g.~the i.MX8MP family,
Jetson Xavier/Orin line) are a third option. Each sensor has its own
receiver and its own DMA pipeline; there is no shared lane to multiplex.
The community threads for the i.MX8MM dual- OV5640 case (NXP Community
posts 1064767 and 1435731) discuss a variety of configurations, ranging
from single-VC with GPIO mux (analogous to our board) to dual-CSI with
fully independent streams. The trade-off is not framed in the literature
as ``software switch cost'' but as ``how many lanes the application of
interest can afford''; our measurement of what the single-lane topology
costs is, to our knowledge, quantitative evidence that was previously
only inferred.

\subsection{1.4 Summary of topology
options}\label{summary-of-topology-options}

\textbf{Table 1.} Comparative positioning of dual-camera topologies,
with per-topology cost of a camera switch.

{\def\LTcaptype{none} % do not increment counter
\begin{longtable}[]{@{}
  >{\raggedright\arraybackslash}p{(\linewidth - 6\tabcolsep) * \real{0.2500}}
  >{\raggedright\arraybackslash}p{(\linewidth - 6\tabcolsep) * \real{0.2500}}
  >{\raggedright\arraybackslash}p{(\linewidth - 6\tabcolsep) * \real{0.2500}}
  >{\raggedright\arraybackslash}p{(\linewidth - 6\tabcolsep) * \real{0.2500}}@{}}
\toprule\noalign{}
\begin{minipage}[b]{\linewidth}\raggedright
Topology
\end{minipage} & \begin{minipage}[b]{\linewidth}\raggedright
Shared CSI?
\end{minipage} & \begin{minipage}[b]{\linewidth}\raggedright
Switch cost per flip
\end{minipage} & \begin{minipage}[b]{\linewidth}\raggedright
Hardware complexity
\end{minipage} \\
\midrule\noalign{}
\endhead
\bottomrule\noalign{}
\endlastfoot
Analogue MUX + single CSI receiver (\textbf{this work}) & yes &
\textasciitilde83 ms (\url{evaluation.md} Table 3) & minimum --- one
analogue switch IC \\
MIPI-bridge with VC split + single CSI receiver & yes, but demultiplexed
in hardware & 0 ms (both live) & moderate --- bridge chip, extra board
area \\
Two independent CSI receivers on SoC & no & 0 ms (software ``select'' is
a buffer choice) & high --- SoC dependency, routing \\
\end{longtable}
}

\begin{center}\rule{0.5\linewidth}{0.5pt}\end{center}

\section{2. VSYNC-synchronised switching and frame-boundary
techniques}\label{vsync-synchronised-switching-and-frame-boundary-techniques}

The pattern of ``change hardware state only during the blanking
interval'' is old in display systems and has well-known analogues in
capture. The camera-capture literature (RidgeRun wiki on Camera Sensor
Basics; Raspberry Pi forum thread 48238 on frame synchronisation; Jetson
developer forum 252751 on CSI-2 sync packet shifts) uses the term
``synchronisation to the frame boundary'' in two distinct ways:

\begin{enumerate}
\def\labelenumi{\arabic{enumi}.}
\tightlist
\item
  \textbf{Sync two cameras to each other} (using FSIN on the OV5640,
  master/slave wiring). This is the approach taken when the system needs
  two \emph{simultaneously captured} frames, typically for stereo or for
  HDR with two different exposures. It does not solve the switching
  problem --- both sensors still stream, and if only one receiver is
  available, the second is lost unless there is a MUX.
\item
  \textbf{Sync the MUX flip to the frame boundary} of the current
  sensor. This is what our flip-on-EOF primitive does. The literature
  for this specific pattern in embedded camera pipelines is sparse ---
  most industrial and NVIDIA Jetson-side systems assume dual receivers
  --- so our contribution is a concrete implementation and a measurement
  of the residual cost.
\end{enumerate}

Framebuffer-replacement techniques in display drivers (double buffering,
page flipping, Vsync-triggered DMA update) are the architectural
ancestor of flip-on-EOF. The analogue is not exact because we are
flipping an input multiplexer rather than swapping an output buffer, but
the principle --- \emph{do the hardware state change during the interval
when no data is flowing} --- is identical.

\begin{center}\rule{0.5\linewidth}{0.5pt}\end{center}

\section{3. Sensor-side effects after a
switch}\label{sensor-side-effects-after-a-switch}

Once the mid-buffer seam is eliminated (§2), the residual visual
artefact observed at \texttt{switch\_drain=1} in our E17 data
(\url{threats_to_validity.md} §2.5) is consistent with known OV5640
post-stream-resume behaviour: AEC/AGC loops take several frames to
converge, and the first one or two frames after a line becomes active
again can be mildly over- or under-exposed relative to the sensor's
final state. The Armbian forum thread (armbian.com topic 4688) on OV5640
color-space configuration, and the ESP32-camera GitHub issue 201 on
frame-rate regressions, both touch on this effect obliquely. We did not
find a published quantitative characterisation of how many frames are
needed for AEC/AGC convergence on this specific sensor; our choice of
\texttt{switch\_drain=2} as the production default is an engineering
conservatism grounded in the observation that three separate sessions
(\href{../experiments/s039_e17_drain_ab/}{s039},
\href{../experiments/s041_e17_eof_check/}{s041}, the user-visual review)
all showed some residual artefact at \texttt{drain=1}. A follow-up study
could instrument the sensor's AEC/AGC status register to tie the
observed pixel artefact to the sensor's own convergence counter.

\begin{center}\rule{0.5\linewidth}{0.5pt}\end{center}

\section{4. ISR-context GPIO and NASA/JPL-style real-time
discipline}\label{isr-context-gpio-and-nasajpl-style-real-time-discipline}

The implementation choice to do the MUX flip inside the CSI EOF ISR is
small in code volume but significant in real-time terms. Three strands
of prior work inform this choice:

\begin{itemize}
\tightlist
\item
  \textbf{AUTOSAR / industrial real-time} --- the principle that ISRs
  should be minimal and defer work to task context is codified in
  AUTOSAR methodology and widely applied in automotive MCU firmware. Our
  ISR does not defer --- it performs the GPIO write directly --- but it
  is still compliant with the ``minimum, bounded, deterministic'' rule
  because the operation is a single atomic-write to \texttt{DR\_SET} /
  \texttt{DR\_CLEAR} (see
  \href{../../../libs/base/gpio.cc}{libs/base/gpio.cc}
  \texttt{GpioSetFromIsr}) with no mutex, no loop, no scheduler
  interaction.
\item
  \textbf{NASA/JPL flight-software rules} --- the discipline referenced
  in \href{../agent/embeded.md}{agent/embeded.md} (bounded loops,
  bounded ISRs, explicit fault containment, error-code taxonomy). Our
  implementation explicitly applies the rules: ISR body is branch-
  and-store only; ratio scheduling is stateless modulo; task-side
  timeout uses delta arithmetic surviving a tick-counter wrap; every
  degraded path has its own error code and counter.
\item
  \textbf{RT-MCU embedded patterns} --- the broader body of MCU-camera
  literature routinely uses GPIO writes in ISRs for pin-state changes
  that must race against hardware timing (e.g.~exposure triggering,
  flash-strobe). What is less common, and is our contribution, is wiring
  the ISR write to a measured MIPI-CSI timing event (end-of-frame)
  rather than to a free-running timer.
\end{itemize}

The end result is that the timing of the MUX flip is no longer
software-controlled in the sense of ``when does the task happen to
run''; it is hardware-event-controlled. That is the transition from
best-effort to deterministic that this work makes concrete for the
dual-sensor case.

\begin{center}\rule{0.5\linewidth}{0.5pt}\end{center}

\section{5. Cross-platform framing}\label{cross-platform-framing}

MCU-class camera-inference pipelines on NXP i.MX RT family (Zephyr
support for the VMU RT1170 board, NXP application notes AN13116,
AN13264, AN5305) focus on single-sensor throughput and power
optimisation; dual-sensor switching is occasionally mentioned as a
feature of the platform but not, to our knowledge, measured to the
per-stage granularity we present. Our E18 head-to-tail benchmark is
portable --- any platform implementing the same \texttt{diag} session
model can run the same three-sweep A/B/C benchmark against any MUX
topology and any sensor pair --- and would be useful as a comparison
point for future platforms claiming lower switch cost.

OpenMV's OV5640 breakout, Adafruit's OV5640 module, and the broader
Arduino-ecosystem OV5640 boards expose the same sensor family but do not
target the dual-sensor use case we focus on. Their configurations are
single-camera and their performance envelopes are therefore not directly
comparable.

\begin{center}\rule{0.5\linewidth}{0.5pt}\end{center}

\section{6. References to gather for the
bibliography}\label{references-to-gather-for-the-bibliography}

Indicative external sources identified during this study; to be folded
into the paper's \texttt{references.bib}. Each is cited in this or the
companion chapters at the point where its content is used.

{\def\LTcaptype{none} % do not increment counter
\begin{longtable}[]{@{}
  >{\raggedright\arraybackslash}p{(\linewidth - 2\tabcolsep) * \real{0.5000}}
  >{\raggedright\arraybackslash}p{(\linewidth - 2\tabcolsep) * \real{0.5000}}@{}}
\toprule\noalign{}
\begin{minipage}[b]{\linewidth}\raggedright
Source
\end{minipage} & \begin{minipage}[b]{\linewidth}\raggedright
Cited for
\end{minipage} \\
\midrule\noalign{}
\endhead
\bottomrule\noalign{}
\endlastfoot
NXP Community post 1064767 --- ``Dual camera(ov5640) in iMx8 M mini'' &
dual-OV5640 configuration on shared CSI with GPIO mux \\
NXP Community post 1992533 --- ``OV5640 and OV7251 with IMX8MM over
MIPI-CSI'' & mixed-sensor dual-camera topology \\
NXP Community post 825091 --- ``two OV5640 mipi csi cameras'' &
single-VC, single-receiver dual-OV5640 \\
NXP AN5305 --- ``MIPI-CSI2 Peripheral on i.MX6 MPUs'' & MIPI-CSI2 packet
structure, FS/LS semantics \\
TI E2E forum 874026 --- ``TS5MP645: MIPI CSI-2 multiplexer between
camera and processor'' & impedance and timing implications of analogue
MUX on a D-PHY \\
RidgeRun wiki --- ``Camera Sensor Basics'' & VSYNC / HSYNC / FS / LS
vocabulary and sync-packet behaviour \\
NXP IMXRT1170 reference manual & CSI peripheral EOF / SOF interrupt
capability \\
OmniVision OV5640 datasheet & 720p@60fps via 2×2 binning; available
modes per resolution \\
Armbian forum topic 4688 --- ``OV5640 on OPI0+: How to reach higher fps
and color spaces'' & anecdotal evidence on OV5640 AEC/AGC post-resume
behaviour \\
Adafruit OV5640 breakout documentation & FSIN slave-sync pin
availability \\
Bewley et al., \emph{SORT}, 2016 & tracking-by-detection association
baseline, referenced also in \url{relatedwork.md} \\
Zhang et al., \emph{ByteTrack}, 2022 & low-score detection recovery in
tracking \\
\end{longtable}
}

The full BibTeX entries live in the existing
\href{references.bib}{\texttt{references.bib}}.
